\section{Estimation details}
\label{app:estimation}

\paragraph{Marginal maximum likelihood.} Writing $\Phi = (\diff, a)$ for the
structural parameters and $p(\skill) = \mathcal{N}(0,1)$ for the population skill
density, the marginal likelihood of rider $\rider$'s outcomes
$\{(\metric_{ri}, y_{ri})\}_i$ is
\begin{equation}
  L_\rider(\Phi) = \int \prod_i
    \sigmoid\!\big(a(\metric_{ri})(\skill - \diff(\metric_{ri}))\big)^{y_{ri}}
    \big(1 - \sigmoid(\cdot)\big)^{1 - y_{ri}}\; p(\skill)\, d\skill .
\end{equation}
The EM algorithm alternates an E-step, which forms the posterior over each rider's
skill on a fixed quadrature grid, and an M-step, which maximises the expected
complete-data log-likelihood over $\Phi$. The full log-likelihood is
$\sum_\rider \log L_\rider(\Phi)$.

\paragraph{Gauss--Hermite quadrature.} The skill integral is evaluated by
Gauss--Hermite quadrature \citep{golub1969}: nodes $\{z_q\}$ and weights
$\{w_q\}$ approximate
\begin{equation*}
  \int f(\skill)\,\mathcal{N}(\skill;0,1)\,d\skill
    \;\approx\; \sum_q w_q\, f(z_q)
\end{equation*}
after the standard change of variables. The node count
trades accuracy against cost; the demonstration configuration uses a fixed grid
adequate for a unimodal skill posterior.

\paragraph{Difficulty representation and anchoring.} The difficulty function is
represented on a metric grid and evaluated by interpolation. Its prior is centred
on the physics-derived expert curve, which anchors the difficulty origin and scale
(Section~\ref{sec:ident}); rare or unobserved conditions shrink toward that
anchor. Rider skills carry the $\mathcal{N}(0,1)$ population prior, which both sets
the location/scale convention and shrinks separated (all-success or all-failure)
riders to finite values.

\paragraph{Determinism.} Given a cohort and a seed, the estimator is
deterministic---no stochastic sampling enters the point estimate---so results are
bit-for-bit reproducible.

\section{Data-generating process and reproducibility}
\label{app:repro}

\paragraph{Generator.} The reference cohort is drawn as in Section~\ref{sec:dgp}:
$80$ riders, $30$ outcomes each; skills $\skill_\rider \sim \mathcal{N}(0,1)$;
conditions uniform on $[3,35]$ knots; ground-truth difficulty
\eqref{eq:dgp-delta} with optimum $\metric^\star = 18$~kn, width $w = 8$, floor
$f = -1$; true discrimination $a_{\text{true}} = 2.0$; cohort seed fixed. Outcomes
are Bernoulli draws from the model~\eqref{eq:model}. These are the ground-truth
inputs to the \emph{generator} and are not provided to the fit.

\paragraph{One-command reproduction.} The results in Table~\ref{tab:recovery} are
produced by
\begin{quote}
  \verb|uv run inverse-suitability latent-demo --out DIR|
\end{quote}
which fits the cohort and writes \texttt{stats.json}. Under
\texttt{reproducibility/}, the paper's results table and inline macros are
rendered from that JSON by a small script (\texttt{gen\_results\_tex.py}), and the
difficulty-atlas figure is rendered from the same package's difficulty and link
functions (\texttt{gen\_figure.py}); \texttt{make repro} runs the whole
chain and rebuilds the PDF. Because the reported numbers are generated from the
code rather than transcribed, they cannot drift from it. The run is deterministic,
completes in under a minute, and uses no proprietary data and no network access.

\paragraph{Recovered versus ground-truth values.} For the avoidance of doubt: the
values in Table~\ref{tab:recovery} are the estimator's \emph{recovered} outputs
(for example, the difficulty minimum recovered at
$\hat\metric^\star = \difficultyMinKnots$~kn). The value $\metric^\star = 18$~kn is
the \emph{ground truth} supplied to the generator. The recovery claim is precisely
that the former lands within $\pm 3$~kn of the latter.
